\documentclass[10pt,a4paper]{amsart}
\usepackage[margin=19mm]{geometry}
\usepackage{amsmath,amssymb,mathtools}
\usepackage[scr=rsfs,cal=euler]{mathalfa}
\usepackage{xfrac}
\usepackage[unicode,hidelinks,bookmarks=false]{hyperref}
\newcommand{\ie}{i.e.\ }
\newcommand{\cf}{cf.\ }
\newcommand{\resp}{respectively\ }
\newcommand{\Subsection}{\subsection}
\providecommand{\nobreakdash}{\nobreak\hbox{-}}
\setlength{\emergencystretch}{2em}
\multlinegap0pt
\usepackage{amssymb}
\usepackage[all]{xy}
\usepackage{tikz}
\newtheorem{lem}{Lemma}
\newcommand{\Nq}{N_\Qo}
\newcommand{\Qo}{\mathbb{Q}}
\title{Toric orbit spaces which are manifolds}
\author{Anton Ayzenberg, Vladimir Gorchakov}
\date{Source: arXiv:2302.02058v1 (CC BY 4.0)}
\begin{document}
\maketitle

\noindent\emph{Source and licence.} Toric orbit spaces which are manifolds. Version arXiv:2302.02058v1 (CC BY 4.0), distributed by arXiv under the Creative Commons Attribution 4.0 International licence, http://creativecommons.org/licenses/by/4.0/, as recorded in the arXiv licence metadata for this exact version and shown on the abstract page https://arxiv.org/abs/2302.02058v1. Full licence notice: \texttt{licenses/arxiv-2302.02058-CC-BY-4.0.txt}.\par\smallskip

\noindent\textit{Editorial scope.} Counted unit: the article's lem lem (source0.tex lines 550--552). The article's complete original proof of it is delivered with it (source0.tex lines 554--556, one \texttt{proof} environment of 3 lines). No mathematics is written, repaired or re-derived: the statement and the proof are the article's own lines, in the article's own notation, and no retained line is substituted or re-flowed.  No further source-local context is needed: every cross-reference of every retained line resolves inside the two delivered spans.  Nothing was omitted from any retained span, so the package carries no omission record.  No retained line contains a citation, so the wrapper carries no bibliography and no bibliography entry.  No retained line contains a figure environment, an image-inclusion command or drawing code, so no image is delivered and the package declares no asset dependency.  Editorial formatting only, in addition: an AMS \texttt{amsart} preamble that restates the article's own declarations and macros used by the retained lines, namely the environments \texttt{lem} on separate counters, together with the standard packages those lines need.  The retained environments print in this document as Lemma 1 in order of appearance, whereas the article prints the same items with the numbering of its own full document.  The complete native source of the article as distributed in the arXiv e-print for this exact version is bundled unchanged as \texttt{sources/136959/source0.tex}.
\begin{lem}
Every representation of complexity one is a Leontief representation.
\end{lem}

\begin{proof}
A representation of complexity one is characterized by $n$ weights $\alpha_1,\ldots,\alpha_n$ in the $(n-1)$-dimensional vector space $\Nq\cong\Qo^{n-1}$. Hence there is a unique (up to multiplier) linear relation $c_1\alpha_1+\cdots+c_n\alpha_n=0$. The weights' subset $\{\alpha_i\mid c_i\neq0\}$ corresponds to a complexity one action in general position, while the remaining weights correspond to an action of complexity $0$.
\end{proof}

\end{document}
